\section{Choi-type maps}

The Choi-type maps in Wolf's Example~3.1 are built from the diagonal projection
$D(X)$ and cyclic shifts $U_{k0}$:
\begin{align}
    T_C(X)
     & =(d-n)D(X)-X+\sum_{k=1}^nD(U_{k0}XU_{k0}^{\dagger}).
    \notag
\end{align}
The formal statements below record the map and the rank-one reduction used in
the standard positivity argument. Positivity is proved, for every $d\ge3$, at
both ends of the range: $n=1$ and $n=d-2$. The positivity assertion for the
middle range $2\le n\le d-3$ and the indecomposability assertion remain open.
The missing general positivity step is a cyclic reciprocal estimate for the
weights in the rank-one reduction; it does not follow merely from equality of
the total numerator and denominator weights. The general estimate is
classical: the case $n=d-2$ is due to Ando, while the case $n=1$ is Theorem~1
of Tanahashi and Tomiyama~\cite[pp.~309--311]{TanahashiTomiyama1988Indecomposable}.
Their Lemmas~2--3 use the same rank-one and product-one coefficient argument
formalized below. The whole range is due to Yamagami~\cite{Yamagami1993Cyclic},
whose proof is a variational argument on the boundary of a projective simplex.

\begin{definition}[Choi-type map]
    \label{def:choi_type_map}
    \lean{Matrix.choiTypeMap}
    \leanok
    On the cyclic index set $\Z/d\Z$, the Choi-type map is
    \begin{align}
        T_C(X)
         & =(d-n)D(X)-X+\sum_{k=1}^nD(U_{k0}XU_{k0}^{\dagger}),
        \notag
    \end{align}
    where $D$ keeps the diagonal of a matrix and $U_{k0}$ is the cyclic shift by
    $k$.
\end{definition}

\begin{lemma}[The Choi-type map on rank-one projectors]
    \label{lem:choi_type_map_rankone}
    \lean{Matrix.choiTypeMap_vecMulVec}
    \leanok
    \uses{def:choi_type_map}
    For a vector $v$ on $\Z/d\Z$,
    \begin{align}
        T_C(\ketbra{v}{v})
         & =
        \diag\!\left((d-n)|v_i|^2+\sum_{k=1}^n |v_{i-k}|^2\right)
        -\ketbra{v}{v}.
        \notag
    \end{align}
\end{lemma}

\begin{proof}
    \leanok
    The diagonal projection of a shifted rank-one projector records the shifted
    squared moduli:
    $D(U_{k0}\ketbra{v}{v}U_{k0}^{\dagger})_{ii}=|v_{i-k}|^2$.
    Substituting this into the definition of $T_C$ gives the displayed diagonal
    matrix minus the original rank-one projector.
\end{proof}

\begin{definition}[The Choi rank-one diagonal weight]
    \label{def:choi_type_rankone_weight}
    \lean{Matrix.choiTypeRankOneWeight}
    \leanok
    For a vector $v$ on $\Z/d\Z$, set
    \begin{align}
        a_i
         & =(d-n)|v_i|^2+\sum_{k=1}^n |v_{i-k}|^2.
        \notag
    \end{align}
\end{definition}

\begin{lemma}[Rank-one complement from a unit bound]
    \label{lem:rankone_complement_from_unit_bound}
    \lean{Matrix.one_sub_vecMulVec_posSemidef_of_sum_normSq_le_one}
    \leanok
    If $p$ is a vector with $\sum_i |p_i|^2\le 1$, then
    $\Id-\ketbra{p}{p}\ge 0$.
\end{lemma}

\begin{proof}
    \leanok
    For every vector $w$,
    $\bra{w}(\Id-\ketbra{p}{p})\ket{w}
        =\|w\|^2-|\braket{p}{w}|^2$.
    Since $\|p\|^2=\sum_i |p_i|^2\le 1$, Cauchy's inequality gives
    $|\braket{p}{w}|^2\le \|p\|^2\|w\|^2\le \|w\|^2$, so the quadratic form is
    non-negative.
\end{proof}

\begin{lemma}[Diagonal rank-one Schur-complement criterion]
    \label{lem:diagonal_rankone_schur_complement}
    \lean{Matrix.diagonal_sub_vecMulVec_posSemidef_of_sum_normSq_div_le_one}
    \leanok
    \uses{lem:rankone_complement_from_unit_bound}
    Let $a_i\ge0$ and suppose $v_i=0$ whenever $a_i=0$. If, with the convention
    that $|v_i|^2/a_i$ is read as $0$ when $a_i=0$,
    $\sum_i |v_i|^2/a_i\le 1$, then
    $\diag(a_i)-\ketbra{v}{v}\ge 0$.
\end{lemma}

\begin{proof}
    \leanok
    Put $p_i=0$ when $a_i=0$ and $p_i=v_i/\sqrt{a_i}$ otherwise. Then
    \begin{align}
        \|p\|^2
         & =\sum_i |p_i|^2
        =\sum_i \frac{|v_i|^2}{a_i}
        \le 1
        \notag
    \end{align}
    under the same zero-term convention.
    Lemma~\ref{lem:rankone_complement_from_unit_bound}
    gives $\Id-\ketbra{p}{p}\ge0$. If $S=\diag(\sqrt{a_i})$, then
    the support condition gives
    $S(\Id-\ketbra{p}{p})S^\dagger=\diag(a_i)-\ketbra{v}{v}$.
    Conjugation by $S$ preserves positive semidefiniteness, which proves the
    claim.
\end{proof}

\begin{lemma}[Rank-one positivity from the cyclic reciprocal bound]
    \label{lem:choi_type_rankone_positive_from_weight_bound}
    \lean{Matrix.choiTypeMap_vecMulVec_posSemidef_of_weight_sum_le_one}
    \leanok
    \uses{def:choi_type_map, lem:choi_type_map_rankone,
        def:choi_type_rankone_weight, lem:diagonal_rankone_schur_complement}
    Assume $n\le d-2$. If, with the convention that the summand is $0$ at
    indices where $a_i=0$, $\sum_i |v_i|^2/a_i\le 1$, where $a_i$ is the weight
    of Definition~\ref{def:choi_type_rankone_weight}, then
    $T_C(\ketbra{v}{v})\ge 0$.
\end{lemma}

\begin{proof}
    \leanok
    The matrix in Lemma~\ref{lem:choi_type_map_rankone} has the form
    $\diag(a_i)-\ketbra{v}{v}$. Each $a_i$ is non-negative. Since
    $n\le d-2$, the coefficient $d-n$ is positive; hence $a_i=0$ forces
    $v_i=0$. The hypotheses of
    Lemma~\ref{lem:diagonal_rankone_schur_complement} therefore apply and give
    the stated positivity.
\end{proof}

\begin{lemma}[The first cyclic reciprocal inequality]
    \label{lem:choi_type_first_cyclic_reciprocal}
    \lean{Matrix.choiType_cyclic_reciprocal_one_of_pos,
        Matrix.choiType_cyclic_reciprocal_one_of_nonneg,
        Matrix.choiTypeRankOneWeight_one,
        Matrix.choiTypeRankOneWeight_reciprocal_sum_one}
    \leanok
    Let $d\ge2$ and let $x_i\ge0$ be indexed by $\Z/d\Z$. Then, with
    zero-denominator summands read as $0$,
    \begin{align}
        \sum_i\frac{x_i}{(d-1)x_i+x_{i-1}}
         & \le1.
        \notag
    \end{align}
    Consequently, for $n=1$ and $x_i=|v_i|^2$, the reciprocal sum of the
    rank-one weights $a_i=(d-1)x_i+x_{i-1}$ is at most $1$.
\end{lemma}

\begin{proof}
    \leanok
    Following Tanahashi--Tomiyama, Lemma~3, first suppose that every $x_i$ is
    positive and set
    $y_i=x_{i-1}/x_i$. Then $\prod_i y_i=1$, and the sum becomes
    $\sum_i 1/(d-1+y_i)$. After clearing its positive common denominator,
    multiply the resulting non-negative difference by the positive factor
    $d-1$. The desired estimate then reduces to
    \begin{align}
        \sum_{k=0}^{d}(k-1)(d-1)^{d-k}e_k(y)
         & \ge0,
        \notag
    \end{align}
    where $e_k(y)$ is the $k$th elementary symmetric sum. The product of the
    $k$-fold monomials occurring in $e_k(y)$ is $1$, so AM--GM gives
    $e_k(y)\ge\binom{d}{k}$. The coefficients are non-negative for $k\ge2$,
    while the $k=0,1$ terms are fixed, and substitution of
    $e_k(y)=\binom{d}{k}$ gives equality by the binomial identity obtained at
    $y_i=1$.

    Tanahashi--Tomiyama assume positive coordinates at this step. To obtain the
    source's unrestricted positive-map statement, suppose instead that some
    $x_i$ vanishes. Then at most $d-1$ summands are nonzero. Every nonzero
    summand is at most $1/(d-1)$ because
    $(d-1)x_i+x_{i-1}\ge(d-1)x_i$, so their sum is again at most $1$.
\end{proof}

\begin{lemma}[Rank-one positivity at the bottom of the range]
    \label{lem:choi_type_rankone_positive_bottom}
    \lean{Matrix.choiTypeMap_vecMulVec_posSemidef_one}
    \leanok
    \uses{lem:choi_type_first_cyclic_reciprocal,
        lem:choi_type_rankone_positive_from_weight_bound}
    For $d\ge3$ and $n=1$, every vector $v$ satisfies
    $T_C(\ketbra{v}{v})\ge0$.
\end{lemma}

\begin{proof}
    \leanok
    Put $x_i=|v_i|^2$. Lemma~\ref{lem:choi_type_first_cyclic_reciprocal}
    bounds $\sum_i x_i/a_i$ by $1$, and
    Lemma~\ref{lem:choi_type_rankone_positive_from_weight_bound} turns this
    bound into rank-one positivity.
\end{proof}

\begin{lemma}[Positivity from rank-one positivity]
    \label{lem:positive_map_from_rankone}
    \lean{Matrix.isPositiveMap_of_forall_vecMulVec_posSemidef}
    \leanok
    A linear map on matrices that sends every rank-one projector
    $\ketbra{w}{w}$ to a positive semidefinite matrix sends every positive
    semidefinite matrix to a positive semidefinite matrix.
\end{lemma}

\begin{proof}
    \leanok
    Let $\Phi$ be the map and let $X\ge0$. Scaling the eigenvectors of $X$ by
    the square roots of its non-negative eigenvalues produces vectors
    $w_\alpha$ with $X=\sum_\alpha \ketbra{w_\alpha}{w_\alpha}$.
    Linearity gives
    $\Phi(X)=\sum_\alpha \Phi(\ketbra{w_\alpha}{w_\alpha})$, each summand is
    positive semidefinite by hypothesis, and a sum of positive semidefinite
    matrices is positive semidefinite.
\end{proof}

\begin{theorem}[Choi-type maps at the bottom of the range are positive]
    \label{thm:choi_type_positive_bottom}
    \lean{Matrix.choiTypeMap_isPositiveMap_one}
    \leanok
    \uses{def:choi_type_map, lem:choi_type_rankone_positive_bottom,
        lem:positive_map_from_rankone}
    For every $d\ge3$ and $n=1$, the Choi-type map $T_C$ sends every positive
    semidefinite matrix to a positive semidefinite matrix. This is the bottom
    slice of the positivity assertion of Wolf's Example~3.1.
\end{theorem}

\begin{proof}
    \leanok
    Combine Lemma~\ref{lem:choi_type_rankone_positive_bottom} with
    Lemma~\ref{lem:positive_map_from_rankone}.
\end{proof}

\begin{lemma}[Permutation reciprocal inequality]
    \label{lem:choi_type_perm_reciprocal}
    \lean{Matrix.perm_reciprocal_sum_le_one}
    \leanok
    Let $x_i\ge0$ be indexed by a finite set, let $T=\sum_j x_j$, and let
    $\sigma$ be a permutation of the index set. Then, with zero-denominator
    summands read as $0$,
    \begin{align}
        \sum_i \frac{x_i}{T+x_i-x_{\sigma(i)}}
         & \le 1.
        \notag
    \end{align}
\end{lemma}

\begin{proof}
    \leanok
    If $T=0$, every summand vanishes. Otherwise write
    $\delta_i=x_i-x_{\sigma(i)}$. Each summand obeys
    \begin{align}
        \frac{x_i}{T+\delta_i}
         & \le\frac{x_iT-x_i\delta_i+\delta_i^2/2}{T^2}.
        \notag
    \end{align}
    For $\sigma(i)=i$ both sides equal $x_i/T$, and for $\sigma(i)\ne i$ the
    difference of the two sides is
    \begin{align}
        \frac{\delta_i^2(T-x_i-x_{\sigma(i)})}
        {2T^2(T+\delta_i)}
         & \ge0,
        \notag
    \end{align}
    by the pair bound $x_i+x_{\sigma(i)}\le T$ and
    $T+\delta_i\ge 2x_i\ge0$; a vanishing denominator forces $x_i=0$ and leaves
    a non-negative right-hand side. Because $\sigma$ permutes the entries,
    $\sum_i x_{\sigma(i)}^2=\sum_i x_i^2$, which gives the identity
    $\sum_i x_i\delta_i=\tfrac12\sum_i\delta_i^2$. Summing the term bounds
    therefore yields
    \begin{align}
        \sum_i\frac{x_i}{T+\delta_i}
         & \le
        \frac{T^2-\sum_i x_i\delta_i+\tfrac12\sum_i\delta_i^2}{T^2}
        =1.
        \notag
    \end{align}
\end{proof}

\begin{lemma}[The Choi weight at the top of the range]
    \label{lem:choi_type_weight_top}
    \lean{Matrix.sum_shift_window_eq_sub_pair,
        Matrix.choiTypeRankOneWeight_sub_two}
    \leanok
    \uses{def:choi_type_rankone_weight}
    Let $d\ge3$ and $n=d-2$. With $x_i=|v_i|^2$ and $T=\sum_j x_j$, the Choi
    rank-one diagonal weight equals
    \begin{align}
        a_i
         & =(d-n)x_i+\sum_{k=1}^{n}x_{i-k}
        =T+x_i-x_{i+1}.
        \notag
    \end{align}
\end{lemma}

\begin{proof}
    \leanok
    For $n=d-2$ the backward shifts $i-1,\ldots,i-(d-2)$ run over every index
    of $\Z/d\Z$ except $i$ and $i+1$, so the shifted sum equals
    $T-x_i-x_{i+1}$ and $a_i=2x_i+T-x_i-x_{i+1}$.
\end{proof}

\begin{lemma}[Rank-one positivity at the top of the range]
    \label{lem:choi_type_rankone_positive_top}
    \lean{Matrix.choiTypeMap_vecMulVec_posSemidef_sub_two}
    \leanok
    \uses{lem:choi_type_rankone_positive_from_weight_bound,
        lem:choi_type_weight_top, lem:choi_type_perm_reciprocal}
    For $d\ge3$ and $n=d-2$, every vector $v$ satisfies
    $T_C(\ketbra{v}{v})\ge0$.
\end{lemma}

\begin{proof}
    \leanok
    By Lemma~\ref{lem:choi_type_weight_top}, the reciprocal weight sum is
    $\sum_i x_i/(T+x_i-x_{i+1})$ with $x_i=|v_i|^2$.
    Lemma~\ref{lem:choi_type_perm_reciprocal}, applied to the cyclic shift
    $\sigma(i)=i+1$, bounds this sum by $1$, and
    Lemma~\ref{lem:choi_type_rankone_positive_from_weight_bound} turns the
    bound into rank-one positivity.
\end{proof}

\begin{theorem}[Choi-type maps at the top of the range are positive]
    \label{thm:choi_type_positive_top}
    \lean{Matrix.choiTypeMap_isPositiveMap_sub_two}
    \leanok
    \uses{def:choi_type_map, lem:choi_type_rankone_positive_top,
        lem:positive_map_from_rankone}
    For every $d\ge3$ and $n=d-2$, the Choi-type map $T_C$ sends every
    positive semidefinite matrix to a positive semidefinite matrix. This is
    the case $n=d-2$ of the positivity assertion of Wolf's Example~3.1 and
    subsumes the case $d=3$, $n=1$.
\end{theorem}

\begin{proof}
    \leanok
    Combine Lemma~\ref{lem:choi_type_rankone_positive_top} with
    Lemma~\ref{lem:positive_map_from_rankone}.
\end{proof}

\section{Transposition}

Transposition is the basic example of a positive map that is not completely
positive. Its positivity is used in the decomposable-map construction below,
while the failure of complete positivity follows from its Choi matrix.

\begin{theorem}[Transposition is positive and trace-preserving]
    \label{thm:transpose_positive_trace_preserving}
    \lean{Matrix.transposeLinearMapComplex_isPositiveMap,
        Matrix.transposeLinearMapComplex_isTracePreservingMap}
    \leanok
    \uses{def:positive_map, def:trace_preserving}
    The matrix transposition map $\theta(X)=X^T$ on $\MN{D}$ is positive
    and trace-preserving.
\end{theorem}

\begin{proof}\leanok
    If $X \ge 0$, write $X=C^\dagger C$. Then
    $X^T=C^T\overline C=\overline C^\dagger\overline C$, so $X^T \ge 0$.
    Trace preservation is the identity $\tr(X^T)=\tr(X)$.
\end{proof}

\begin{theorem}[Transposition is not completely positive]
    \label{thm:transpose_not_completely_positive}
    \lean{ChoiJamiolkowski.transposeLinearMapComplex_not_isCPMap}
    \leanok
    \uses{def:cp_map, def:flip_operator, def:fin_two_antisymm_vec,
        thm:choi_transpose_flip, thm:cp_iff_choi_posSemidef}
    If $D \ge 2$, the matrix transposition map $\theta(X)=X^T$ on $\MN{D}$
    is not completely positive.
\end{theorem}

\begin{proof}\leanok
    \uses{def:flip_operator, def:fin_two_antisymm_vec, thm:choi_transpose_flip,
        thm:cp_iff_choi_posSemidef}
    If $\theta$ were completely positive, its Choi matrix would be positive
    semidefinite. By Theorem~\ref{thm:choi_transpose_flip}, this Choi matrix is
    $D^{-1}F$. For the antisymmetric vector
    $v=\ket{01}-\ket{10}$,
    \begin{align}
        \bra{v}\,D^{-1}F\,\ket{v}
         & =-\frac{2}{D}<0,
        \notag
    \end{align}
    contradicting positive semidefiniteness.
\end{proof}

\section{Decomposable positive maps}

The Choi-type maps are stated in Wolf's Example~3.1 as positive and
indecomposable. The formal language for the second adjective is the following
standard one: a positive map is decomposable when it is the sum of a completely
positive map and a completely copositive map.

\begin{definition}[Decomposable and indecomposable positive maps]
    \label{def:decomposable_positive_map}
    \lean{IsCompletelyCopositiveMap, IsDecomposablePositiveMap,
        IsIndecomposablePositiveMap}
    \leanok
    A map $\Phi\colon\MN{D}\to\MN{D}$ is completely copositive if
    $\Phi\circ\theta$ is completely positive, where $\theta$ is transposition.
    It is decomposable if
    $\Phi=\Phi_{\mathrm{cp}}+\Phi_{\mathrm{ccp}}$ with
    $\Phi_{\mathrm{cp}}$ completely positive and
    $\Phi_{\mathrm{ccp}}$ completely copositive. It is indecomposable if it is
    positive and not decomposable.
\end{definition}

\begin{lemma}[Completely copositive maps are positive]
    \label{lem:completely_copositive_map_positive}
    \lean{IsCompletelyCopositiveMap.isPositiveMap}
    \leanok
    \uses{def:decomposable_positive_map, thm:transpose_positive_trace_preserving,
        thm:cp_is_positive}
    If $\Phi$ is completely copositive, then $\Phi$ is positive.
\end{lemma}

\begin{proof}
    \leanok
    Let $X\ge0$. Transposition preserves positivity, so $\theta(X)\ge0$.
    Since $\Phi\circ\theta$ is completely positive,
    $(\Phi\circ\theta)(\theta(X))\ge0$. Because $\theta^2=\id$, this is exactly
    $\Phi(X)\ge0$.
\end{proof}

\begin{lemma}[Decomposable maps are positive]
    \label{lem:decomposable_map_positive}
    \lean{IsDecomposablePositiveMap.isPositiveMap}
    \leanok
    \uses{def:decomposable_positive_map, thm:cp_is_positive,
        lem:completely_copositive_map_positive}
    If $\Phi$ is decomposable, then $\Phi$ is positive.
\end{lemma}

\begin{proof}
    \leanok
    Write $\Phi=\Phi_{\mathrm{cp}}+\Phi_{\mathrm{ccp}}$, with
    $\Phi_{\mathrm{cp}}$ completely positive and
    $\Phi_{\mathrm{ccp}}$ completely copositive. For $X\ge0$, the two preceding
    positivity results give $\Phi_{\mathrm{cp}}(X)\ge0$ and
    $\Phi_{\mathrm{ccp}}(X)\ge0$. Their sum is positive, hence
    $\Phi(X)\ge0$.
\end{proof}

\begin{theorem}[Decomposable maps preserve positivity on PPT states]
    \label{thm:decomposable_map_tensor_id_ppt}
    \lean{IsDecomposablePositiveMap.tensorMapId_posSemidef_of_isPPT}
    \leanok
    \uses{def:decomposable_positive_map, def:is_ppt}
    If $\Phi\colon\MN{D}\to\MN{D}$ is decomposable and $\rho$ is a positive
    semidefinite matrix on $\C^{D}\otimes\C^{D}$ with the PPT property, then
    $(\Phi\otimes\id)(\rho)\ge0$.
\end{theorem}

\begin{proof}
    \leanok
    \uses{def:decomposable_positive_map, def:is_ppt,
        thm:partial_transpose_left_involution}
    Write $\Phi=\Phi_{\mathrm{cp}}+\Phi_{\mathrm{ccp}}$ with
    $\Phi_{\mathrm{cp}}$ completely positive and $\Phi_{\mathrm{ccp}}$
    completely copositive. Complete positivity implies positivity of every
    finite ampliation, so
    $(\Phi_{\mathrm{cp}}\otimes\id)(\rho)\ge0$. For the second summand,
    ampliating $\Phi_{\mathrm{ccp}}$ after transposition equals ampliating $\Phi_{\mathrm{ccp}}$ after the first-factor partial transpose
    of the input,
    \begin{align}
        (\Phi_{\mathrm{ccp}}\otimes\id)(\rho)
         & =((\Phi_{\mathrm{ccp}}\circ\theta)\otimes\id)(\rho^{T_{1}}),
        \notag
    \end{align}
    because both sides evaluate $\Phi_{\mathrm{ccp}}$ on the transposed
    $(i_{2},j_{2})$-slice of $\rho$. Since $\Phi_{\mathrm{ccp}}\circ\theta$ is
    completely positive and $\rho^{T_{1}}\ge0$ is the PPT hypothesis, the
    right-hand side is positive semidefinite. The sum of the two positive
    images is positive.
\end{proof}

\begin{theorem}[A PPT state detected by an ampliation witnesses indecomposability]
    \label{thm:indecomposable_of_ppt_witness}
    \lean{not_isDecomposablePositiveMap_of_isPPT_not_tensorMapId_posSemidef}
    \leanok
    \uses{def:decomposable_positive_map, def:is_ppt,
        thm:decomposable_map_tensor_id_ppt}
    Let $\Phi\colon\MN{D}\to\MN{D}$ be a linear map. If some positive
    semidefinite matrix $\rho$ on $\C^{D}\otimes\C^{D}$ with the PPT property
    has $(\Phi\otimes\id)(\rho)$ not positive semidefinite, then $\Phi$ is not
    decomposable.
\end{theorem}

\begin{proof}
    \leanok
    This is the contrapositive of
    Theorem~\ref{thm:decomposable_map_tensor_id_ppt}.
\end{proof}
